
\SGAProof{Proof of \SGARefLink{sga3:xvii:lemma:2:3:bis}{Lemma 2.3 bis}}
Indeed, the group of \(k'\)-automorphisms of the algebraic group
\((\mathbb G_a)_{k'}\) is the group of nonzero homotheties
\((k')^\times\) (\textit{Bible}, Exposé 9, \SGARefLink{sga3:xvii:lemma:3:1}{Lemma 1}), and the Galois
cohomology group \(H^1(k'/k,\mathbb G_m)\) is zero (we may suppose that
\(k'\) is a Galois extension of \(k\)), by Hilbert's Theorem 90.
\SGARefLink{sga3:xvii:lemma:2:3:bis}{Lemma 2.3 bis} then follows from the classification of the \(k\)-forms
of an algebraic group (cf.\ J.-P. Serre, \textit{Cohomologie
galoisienne}, Chapter III, \SGARefLink{sga3:iii:subsection:1-3}{1.3}).

This completes the proof of \SGARefLink{sga3:xvii:proposition:2:2}{2.2}.

\SGAStatement{Proposition 2.4}{Let \(k\) be a field, let \(M\) be a
\(k\)-group of multiplicative type (Exposé VIII) and of finite type,
and let \(U\) be a unipotent \(k\)-algebraic group.  Then:
\begin{enumerate}
\item[\SGAInlineRefTarget{sga3:xvii:proposition:2-4:item:i}(i)]\SGAInlineRefTarget{sga3:xvii:proposition:2:4:item:i} \(\underline{\operatorname{Hom}}_{k\text{-gr}}(M,U)
  =\underline e\)
  (a fortiori \(\operatorname{Hom}_{k\text{-gr}}(M,U)=e\));
\item[(ii)] \(\operatorname{Hom}_{k\text{-gr}}(U,M)=e\).
\end{enumerate}}
\SGARefTarget{sga3:xvii:proposition:2:4}

To prove (i), we must establish that, for every prescheme \(S\) over
\(k\),
\[
  \operatorname{Hom}_{S\text{-gr}}(M_S,U_S)=e.
\]
This follows from the next lemma.

\SGAStatement{Lemma 2.5}{Let \(S\) be a prescheme, let \(M\) be an
\(S\)-group of multiplicative type and of finite type over \(S\), and
let \(U\) be an \(S\)-group prescheme, of finite presentation over
\(S\), whose fibers are unipotent.  Then
\(\operatorname{Hom}_{S\text{-gr}}(M,U)=e\).}
\SGARefTarget{sga3:xvii:lemma:2:5}

\SGAProof{Proof}
Under these hypotheses, for an \(S\)-group morphism \(u:M\to U\) to
be the unit homomorphism, it suffices that the restriction of \(u\) to
the fibers of \(M\) over the points of \(S\) be the zero morphism
(\SGARefLink{sga3:ix:corollary:5:2}{Exposé IX, 5.2}).  We are therefore reduced to the case in which \(S\)
is the spectrum of a field, which we may moreover suppose to be
algebraically closed.  In view of \SGARefLink{sga3:xvii:definition:1:1}{Definition 1.1}, it is enough to take
\(U=\mathbb G_a\), in which case the assertion has already been proved
(\SGARefLink{sga3:xii:4.4.1}{Exposé XII, 4.4.1}).

We now prove Proposition 2.4(ii).  Let \(u:U\to M\) be a \(k\)-group
morphism.  The image \(u(U)\) is represented by an algebraic subgroup
\(U''\) of \(M\) (Exposé VI\textsubscript{B}, \SGARefLink{cor:sga3-VIB-5.4}{5.4}).\footnote{Editor's
note: verify this reference.}  The group \(U''\) is unipotent as a
quotient of a unipotent group (2.2(iii)), and is of multiplicative type
as a subgroup of a group of multiplicative type (cf.\ \textit{Bible},
Exposé 4, Theorem 2, \SGARefLink{sga3:xvii:corollary:C:5:1}{Corollary 1}, or \SGARefLink{sga3:ix:theorem:6:8}{Exposé IX, 6.8}).  Thus \(U''\) is
the unit group by \SGARefLink{sga3:xvii:proposition:2-4:item:i}{Proposition 2.4(i)}.

\SGAStatement{\SGAInlineRefTarget{sga3:xvii:remark:2-6}Remark 2.6}{Keeping the notation of \SGARefLink{sga3:xvii:proposition:2:4}{Proposition 2.4}, it
is no longer true in general that the functor
\(\underline{\operatorname{Hom}}_{k\text{-gr}}(U,M)\) is equal to
\(\underline e\).  Indeed, let \(S\) be a prescheme such that
\(\Gamma(S,\mathcal O_S)\) contains a nonzero element \(\varepsilon\)
with \(\varepsilon^2=0\) (for example, the spectrum of the algebra of
dual numbers over a ring \(A\)).  For every \(S'\) over \(S\), the map
\[
  u\longmapsto 1+\varepsilon_{S'}u
\]
defines, functorially in \(S'\), a homomorphism from the additive group
\(\Gamma(S',\mathcal O_{S'})\) to the multiplicative group
\(\Gamma(S',\mathcal O_{S'}^\times)\).  It therefore defines an
\(S\)-group morphism
\[
  (\mathbb G_a)_S\longrightarrow(\mathbb G_m)_S,
\]
and, since \(\varepsilon\ne0\), this morphism is not zero.}
\SGARefTarget{sga3:xvii:remark:2:6}

Recall (Exposé VII\textsubscript{A}, \(\S3\)) that when \(G\) is a
commutative \(S\)-group prescheme, finite and flat over \(S\), Cartier
duality is available and \(G\) is reflexive in the sense of Exposé
VIII, \(\S1\).  More precisely, the functor
\[
  G\longmapsto
  \underline{\operatorname{Hom}}_{S\text{-gr}}(G,\mathbb G_m)
\]
is represented
